\section{Evaluation and maintenance-decision framework}\label{sec:framework}
\subsection{Data, evaluation unit and labels}
C-MAPSS (FD001--FD004)\cite{saxena2008} provides run-to-failure training trajectories and test trajectories truncated before failure, with the terminal RUL supplied separately. FD001 and FD003 have one operating condition and 14 sensor channels; FD002 and FD004 have six conditions, 15 sensor channels and three operating-setting channels. The model input is a window of 30 cycles. Models predict $\min(\mathrm{RUL}_{t+1},125)$, the clipped RUL at the cycle after the window, and test engines are scored against $\min(\mathrm{RUL}-1,125)$ from the official terminal label, so training, calibration, checkpoint selection and testing share one target instant. Every decision quantity below (true RUL, the comparison with the lead time $L$, the clip) refers to that instant. We found no mismatch in the code; scoring against the current-cycle label changes the headline decision rates by at most \JLabelHeadline{} points (\JLabelMax{} over all mechanisms, conditions and lead times).

The evaluation unit is the last window of each test engine. Engines with fewer than 30 recorded cycles cannot form a window and are excluded; the evaluated engines number \JEvalEng{} for FD001--FD004, against \JFileEng{} in the test files. All rates use the evaluated engines as denominator.


\begin{figure}[!t]
\centering
\includegraphics[width=\textwidth]{figures/fig1.pdf}
\caption{Evaluation protocol from clean training-time engine partitions to a single maintenance decision per test engine. The scheme describes the stress test performed here; it is not a deployed monitoring system.}
\label{fig:flow}
\end{figure}

\subsection{Models, calibration and leakage control}
The LSTM has two layers of 64 units with dropout 0.2. The Transformer applies two encoder layers ($d=64$, four heads) with average pooling. Both were trained with Adam for 150 epochs with hyperparameters fixed before the protocol was adopted and not retuned. Five interval mechanisms are built on them: a fixed-variance head (MSE training, one global residual scale), a heteroscedastic Gaussian head (mean and log-scale outputs\cite{nix1994}), MC Dropout ($T=50$, predictive variance including the residual term\cite{gal2016,kendall2017}), a deep ensemble (the five heteroscedastic seeds combined by moment matching\cite{lakshminarayanan2017}) and split conformal prediction with a normalised score (CP-norm\cite{barber2023}).

Figure~\ref{fig:flow} summarises the protocol. Each of five seeds draws its own engine-level partition of the official training file: 60\% of engines fit the weights and the min--max normaliser, 20\% select the checkpoint, and 20\% calibrate the conformal quantile and, for the fixed-variance and MC Dropout mechanisms, the residual scale. The official test file is used once, after all quantities are frozen. Taking the residual scale from the calibration engines rather than from training residuals matters: the calibration-set scale is larger by \NFairRatio\% on average. The protocol repairs an earlier pipeline in which the checkpoint had been selected on the test set. A controlled $2\times2$ retraining with the fitting engines held fixed (Supplement S1) shows that test-set checkpoint selection improves RMSE by up to \NLeakRMSEmax\% and the NASA score by up to \NLeakScoremax\% on FD001, whereas the normaliser scope moves RMSE by at most \NNormAbsMaxRMSE{} cycles. The intervals analysed here are therefore not products of test-set selection; the models are not state of the art.

\emph{Predicted-mean processing.} The analyses do not clip the predicted mean identically (Supplement S2). The v1.8 clean benchmark clipped every mean to $[0,125]$; the stress-test and decision analyses leave the heteroscedastic mean (also used by CP-norm) and ensemble members unclipped and clip the other means. The clean baseline here is therefore recomputed through the stress-test code path from the released checkpoints (Table~\ref{tab:clean}); the clipped values stay in S2, labelled v1.8 processing. Clipping the heteroscedastic means changes none of the NLL or CP-norm trigger rates in \JClipRateN{} stored cells and the ensemble rates by at most \JClipRateMaxEns{} points, but it changes coverage under a large mean shift by up to \JClipPICPmaxDrift{} (FD002 drift: \JClipDriftTwoU{} unclipped, \JClipDriftTwoA{} clipped). Coverage under shifted inputs therefore depends on this processing; trigger-error direction and size do not.

\subsection{Sensor degradation}
Perturbations act on the raw readings of the sensor channels of the final window before normalisation, without retraining. Operating-setting channels are never perturbed (a diagnostic that perturbs them is summarised in Section~\ref{sec:sens}). Levels are $\{0.1,0.5,1,2,5\}\%$ of each channel's range over the official training file, a fixed fraction of range and not an instrument full scale. \emph{Noise} is additive Gaussian. \emph{Bias} is a constant per-channel offset with a random sign per channel. \emph{Drift} is a linear ramp from zero to the nominal level at the last step of the window, re-zeroed per window, with the same positive sign on every sensor channel. \emph{Gain} multiplies the raw reading by $1\pm k$ with a random sign per channel, so its level is a percentage of the reading, not of the range. Nominal levels of different families are therefore not comparable. Results are related instead through the measured excursion $f_{\mathrm{oob}}$, the fraction of normalised sensor inputs of the terminal windows outside $[-1,1]$, averaged over the same models and realisations as the coverage. Noise, bias and gain use five seeded realisations; drift is deterministic.

Because drift is sign-coherent and bias is not, a drift-versus-bias contrast confounds temporal structure with cross-channel sign coherence. Three controls separate them (Section~\ref{sec:coherent}): a constant $+5\%$ offset on every sensor channel, a constant $-5\%$ offset, and the ramp rerun as a check. Further controls from the earlier analysis (reversed, shuffled, single-channel and continuous-trajectory drift) are summarised in Section~\ref{sec:controls}.

\subsection{Maintenance rule and error accounting}
For each mechanism the rule computes a one-sided 95\% lower bound $b$ on RUL. For the Gaussian mechanisms $b=\hat\mu-1.645\hat\sigma$, which is also the lower limit of their central 90\% interval; coverage (PICP) refers to that central interval. For CP-norm, $b=\hat\mu-q\hat\sigma$ with $q$ the conformal quantile of the \emph{signed} normalised score at level 0.95, recalibrated on the calibration engines, because the central 90\% quantile of the absolute score is not a one-sided 95\% quantile. Maintenance is triggered when $b\le L$, with $L\in\{10,20,30\}$ cycles.

Let $y$ be the true RUL at the evaluated instant. An \emph{unrecognised maintenance need} is $y\le L$ with no trigger; its rate is $r_{\mathrm{late}}$. A \emph{premature trigger} is a trigger with $y>L+20$; its rate is $r_{\mathrm{early}}$. Other states (a trigger with $y\le L$, a trigger with $L<y\le L+20$, no trigger with $y>L$) carry no cost. Both rates use all evaluated engines as denominator; the conditional unrecognised rate uses only engines with $y\le L$. The cost is $c_{\mathrm{late}}r_{\mathrm{late}}+r_{\mathrm{early}}$ with $c_{\mathrm{late}}\in\{5,20,100\}$, a given binary loss and not a maintenance-cost model.

This is a terminal-instant stress test: one decision per engine at one state. It does not simulate the time of first maintenance, availability, safety events or lifetime cost. Where we quote the remaining life of prematurely triggered engines, it is the true remaining life at the evaluated instant (unclipped, official RUL minus one), and the excess over $L$ is what would be forgone if maintenance followed immediately; it is not an estimate of total life lost.

\subsection{Statistical units and uncertainty}
The resampling unit is the engine. Each rate is the mean over engines of the per-engine mean over the five models (and over five realisations for noise and bias). Intervals are percentile intervals from 2000 bootstrap resamples of engines, applied jointly to all mechanisms, conditions and lead times so that comparisons stay paired. They are conditional on the five trained models, whose initialisation and partition variability are confounded; the model-to-model range of each rate is reported separately. Five seeds times the engines is not an independent sample size, and windows are never treated as engines. Rank reversals across cost ratios are counted descriptively; the bootstrap probability that the top mechanism differs is a stability indicator, not a significance test, with no multiplicity correction over the 96 settings. The engine-level analyses of Sections~\ref{sec:coherent}--\ref{sec:cost} were run once on the released checkpoints, are exploratory, and reproduce the released aggregate decision rates (\JAgreeN{} values, maximum absolute difference \JAgreeMax).
